\section{Maxwell-Gleichungen}
\begin{table}[H]
	\centering
	\begin{tabularx}{\textwidth}{l | X | X}
		\makecell[l]{\textbf{Gaußsches Gesetz}} & 
			\(\boxed{\opdiv \vec{D} = \nabla \cdot \vec{D} = \rho}\) \par \vspace*{3pt} \footnotesize 
				Das elektrische Feld ist ein Quellenfeld. Die Ladung \(Q\) bzw. die Ladungsdichte \(\rho\) sind Quellen des elektrischen Feldes. & 
			\(\boxed{\oiint_{\partial V} \vec{D} \cdot \mathrm{d} \vec{a} = \iiint_V \rho \mathrm{d}V = Q}\) \par \vspace*{3pt} \footnotesize
				Der (elektrische) Fluss durch die geschlossene Oberfläche \(\partial V\) eines Volumens \(V\) ist gleich der elektrischen Ladung in seinem Inneren. \\
		\makecell[l]{\textbf{Gaußsches Gesetz} \\ \textbf{für Magnetismus}} & 
			\(\boxed{\opdiv \vec{B} = \nabla \cdot \vec{B} = 0}\) \par \vspace*{3pt} \footnotesize 
				Das magnetische Feld ist quellenfrei. Es gibt keine magnetischen Monopole. & 
			\(\boxed{\oiint_{\partial V} \vec{B} \cdot \mathrm{d} \vec{a} = 0}\) \par \vspace*{3pt} \footnotesize 
				Der mag. Fluss durch die geschlossene Oberfläche \(\partial V\) eines Volumens \(V\) entspricht der mag. Ladung in seinem Inneren, nämlich 0, da es keine magnetischen Monopole gibt. \\
		\makecell[l]{\textbf{Induktionsgesetz}} & 
			\(\boxed{\oprot \vec{E} = \nabla \times \vec{E} = -\frac{\partial \vec{B}}{\partial t}}\) \par \vspace*{3pt} \footnotesize 
				Jede zeitliche Änderung eines Magnetfeldes bewirkt ein elektrisches Wirbelfeld. & 
			\(\boxed{\oint_{\partial A} \vec{E} \cdot \mathrm{d} \vec{s} = -\iint_A \frac{\partial \vec{B}}{\partial t} \cdot \mathrm{d} \vec{a} = -\frac{\partial \Phi_{eingesch.}}{\partial t}}\)  \par \vspace*{3pt} \footnotesize 
				Die induzierte Umlaufspannung bzgl. der Randkurve \(\partial A\) einer Fläche \(A\) ist gleich der negativen zeitlichen Änderung des magnetischen Flusses durch diese Fläche. \\
		\makecell[l]{\textbf{Amperesches Gesetz} \\ \textbf{Durchflutungsgesetz}} & 
			\(\boxed{\oprot \vec{H} = \nabla \times \vec{H} = \vec{j} + \frac{\partial \vec{D}}{\partial t}}\) \par \vspace*{3pt} \footnotesize 
				Jeder Strom und jede zeitliche Änderung des elektrischen Feldes (Verschiebungsstrom) bewirkt ein magnetisches Wirbelfeld & 
			\(\boxed{\oint_{\partial A} \vec{H} \cdot \mathrm{d} \vec{s} = \iint_A \left( \vec{j} + \frac{\partial \vec{D}}{\partial t} \right) \cdot \mathrm{d} \vec{a}}\)  \par \vspace*{3pt} \footnotesize 
				Die mag. Umlaufspannung bzgl. der Randkurve \(\partial A\) einer Fläche \(A\) entspricht dem von dieser Fläche eingeschlossenen Strom (inkl. Verschiebungsstrom).
	\end{tabularx}
\end{table}

\subsection{Materialgleichungen}
	\begin{table}[H]
		\centering
		\begin{tabular}{l  l  l}
			\textbf{Differentielles ohmsches Gesetz} & \textbf{Materialgleichung} & \textbf{Materialgleichung} \\
			\(\boxed{\vec{J} = \kappa \cdot \vec{E} = \left[\frac{\mathrm{A}}{\mathrm{m}^2}\right]}\) & 
			\(\boxed{\vec{D} = \varepsilon \cdot \vec{E} = \left[\frac{\mathrm{C}}{\mathrm{m}^2}\right]}\) & 
			\(\boxed{\vec{B} = \mu \cdot \vec{H} = \left[\mathrm{T}\right]}\)
		\end{tabular}
	\end{table}

\subsection{Feldstärkekomponenten einer ebenen Welle}
	Bei Ausbreitung in $z$-Richtung gibt es keine Amplitudenabhängigkeit von $x, y$
	d.h. $\frac{\partial \ldots}{\partial x}=\frac{\partial \ldots}{\partial y}=0$ \\
	Ableitung einer harmonischen Größe nach der Zeit ergibt $\frac{\partial \ldots}{\partial t}=\mathrm{j} \omega \ldots$ 

	damit ergibt sich aus den Maxwell'schen Gleichungen:
	{\footnotesize
	$$
		\begin{gathered}
			\boxed{\operatorname{rot} \underline{\vec{E}}=-\mathrm{j} \omega \mu \underline{\vec{H}}} \qquad \boxed{\operatorname{rot} \underline{\vec{H}}=\mathrm{j} \omega \varepsilon \underline{\vec{E}}} \\
			\operatorname{rot} \underline{\vec{E}}=\left|\begin{array}{ccc}
				\vec{e}_x                   & \vec{e}_y                   & \vec{e}_z                   \\
				\frac{\partial}{\partial x} & \frac{\partial}{\partial y} & \frac{\partial}{\partial z} \\
				\underline{E}_x             & \underline{E}_y             & \underline{E}_z
			\end{array}\right|=\left(\begin{array}{c}
					\frac{\partial \underline{E}_z}{\partial y}-\frac{\partial \underline{E}_y}{\partial z} \\
					\frac{\partial \underline{E}_x}{\partial z}-\frac{\partial \underline{E}_z}{\partial x} \\
					\frac{\partial \underline{E}_y}{\partial x}-\frac{\partial \underline{E}_x}{\partial y}
				\end{array}\right)=\left(\begin{array}{l}
					0-\frac{\partial \underline{E}_y}{\partial z} \\
					\frac{\partial \underline{E}_x}{\partial z}-0 \\
					0-0
				\end{array}\right)=-\mathrm{j} \omega \mu\left(\begin{array}{l}
					\underline{H}_x \\
					\underline{H}_y \\
					\underline{H}_z
				\end{array}\right) \\
			\operatorname{rot} \underline{\vec{H}}=\left|\begin{array}{ccc}
				\vec{e}_x                   & \vec{e}_y                   & \vec{e}_z                   \\
				\frac{\partial}{\partial x} & \frac{\partial}{\partial y} & \frac{\partial}{\partial z} \\
				\underline{H}_x             & \underline{H}_y             & \underline{H}_z
			\end{array}\right|=\left(\begin{array}{c}
					\frac{\partial \underline{H_z}}{\partial y}-\frac{\partial \underline{H}_y}{\partial z} \\
					\frac{\partial \underline{H}_x}{\partial z}-\frac{\partial \underline{H}_z}{\partial x} \\
					\frac{\partial \underline{H}_y}{\partial x}-\frac{\partial \underline{H}_x}{\partial y}
				\end{array}\right)=\left(\begin{array}{c}
					0-\frac{\partial \underline{H}_y}{\partial z} \\
					\frac{\partial \underline{H}_x}{\partial z}-0 \\
					0-0
				\end{array}\right)=\mathrm{j} \omega \varepsilon\left(\begin{array}{l}
					\underline{E}_x \\
					\underline{E}_y \\
					\underline{E}_z
				\end{array}\right)
		\end{gathered}
	$$
	}

\subsection{Integralsätze}
	\begin{description}
		\setlength{\itemsep}{1pt}
		\item Fundamentalsatz der Analysis
		\item Gauß: Vektorfeld das aus Oberfläche von Volumen strömt muss aus Quelle in Volumen
		\item Stokes: innere Wirbel kompensieren sich $\rightarrow$ nur den Rand betrachten.
	\end{description}
	\begin{align*}
		\int_{a}^b \opgrad F \cdot d \vec{s}     & = F(b) - F(a)                                  \\
		\iiint_V \opdiv \vec{A} \cdot dV         & = \oiint_{ \partial V} \vec{A} \cdot d \vec{a} \\
		\iint_{A} \oprot \vec{A} \cdot d \vec{a} & = \oint_{ \partial A} \vec{A} \cdot d \vec{r}
	\end{align*}


